% Nonarchival additions. Loaded and checked by minds_paper.bibliography.
\begin{thebibliography}{99}
\bibitem{lisiecki2005} L. E. Lisiecki and M. E. Raymo. A Pliocene--Pleistocene stack of 57 globally distributed benthic $\delta^{18}$O records. \emph{Paleoceanography} 20, PA1003, 2005. \doi{10.1029/2004PA001071}.
\bibitem{clark2006} P. U. Clark, D. Archer, D. Pollard, et al. The middle Pleistocene transition: characteristics, mechanisms, and implications for long-term changes in atmospheric pCO$_2$. \emph{Quaternary Science Reviews} 25, 3150--3184, 2006. \doi{10.1016/j.quascirev.2006.07.008}.
\bibitem{margari2023} V. Margari, D. A. Hodell, S. A. Parfitt, et al. Extreme glacial cooling likely led to hominin depopulation of Europe in the Early Pleistocene. \emph{Science} 381, 693--699, 2023. \doi{10.1126/science.adf4445}.
\bibitem{hu2023} W. Hu, Z. Hao, P. Du, et al. Genomic inference of a severe human bottleneck during the Early to Middle Pleistocene transition. \emph{Science} 381, 979--984, 2023. \doi{10.1126/science.abq7487}.
\bibitem{muttoni2024} G. Muttoni and D. V. Kent. Hominin population bottleneck coincided with migration from Africa during the Early Pleistocene ice age transition. \emph{PNAS} 121, e2318903121, 2024. \doi{10.1073/pnas.2318903121}.
\bibitem{deng2025} Y. Deng, R. Nielsen, and Y. S. Song. A previously reported bottleneck in human ancestry 900 kya is likely a statistical artifact. \emph{Genetics} 229, iyae192, 2025 (online December 2024). \doi{10.1093/genetics/iyae192}.
\bibitem{cousins2025} T. Cousins and A. Durvasula. Insufficient evidence for a severe bottleneck in humans during the Early to Middle Pleistocene transition. \emph{Molecular Biology and Evolution} 42, msaf041, 2025. \doi{10.1093/molbev/msaf041}.
\bibitem{zhou2026} S. Zhou, Y. Zhou, F. Di Vincenzo, et al. Severe bottleneck of ancient Homo populations: insights from computational modeling and relevant fossil evidence. \emph{Molecular Biology and Evolution} 43, msag157, July 2026. \doi{10.1093/molbev/msag157}.
\bibitem{sieh2023} K. Sieh, D. Schonwalder Angel, J. Herrin, et al. Proximal ejecta of the Bolaven extraterrestrial impact, southern Laos. \emph{PNAS} 120, e2310351120, 2023. \doi{10.1073/pnas.2310351120}.
\bibitem{mizera2022} J. Mizera. Quest for the Australasian impact crater: failings of the candidate location at the Bolaven Plateau, Southern Laos. \emph{Meteoritics \& Planetary Science} 57, 1973--1986, 2022. \doi{10.1111/maps.13912}.
\end{thebibliography}
